\section{Model}

\newcommand{\vect}[2]{\mathbf{#1}_{#2}}
\newcommand{\norm}[1]{\left\lVert #1 \right\rVert^2}
\newcommand{\trans}[1]{#1^\top}
\newcommand{\Gam}[3]{\mathrm{Gamma} \left( #1 \;\middle|\; #2, #3\right)}
\newcommand{\Nor}[3]{\mathcal{N} \left( #1 \;\middle|\; #2, #3 \right)}
\newcommand{\avgz}{\overline{z}}
\newcommand{\avgzz}{\overline{zz}}

\begin{figure}[h]
    \centering
    \begin{tikzpicture}

% Nodes
        \node[latent] (mu) {$\mu$};
        \node[latent, right=of mu] (tau) {$\tau$};
        \node[latent, below=of mu] (z) {$z_i$};
        \node[obs, below=of z] (X) {$X_{ij}$};
        \node[latent, right=of X] (W) {$\vect{w}{j}$};

% Edges
        \edge {mu,tau} {z};
        \edge {z,W} {X};

% Plates
        \plate {plate1} {(z)(X)} {$n$};
        \plate {plate2} {(X)(W)} {$m$};

    \end{tikzpicture}
    \caption{Bayesian network representing the latent regression model}\label{fig:regression_bayes_net}
\end{figure}


Let $\vect{z}{} = \{z_1, \dots, z_n\}$ be a vector of independent and identically distributed Gaussian latent variables:
\begin{equation*}
    p(\vect{z}{}) = \prod_{i=1}^n \Nor{z_i}{\mu}{\tau^{-1}}.
\end{equation*}

Let $\{\phi_1(\cdot), \dots, \phi_k(\cdot)\}$ denote a collection of basis functions.
We define the vector-valued feature map
\[
    \phi(z_i) = (\phi_1(z_i), \dots, \phi_k(z_i))^\top \in \mathbb{R}^k.
\]
The design matrix $Z \in \mathbb{R}^{n \times k}$ is then defined by stacking these vectors row-wise, so that
\[
    Z_{i\ell} = \phi_\ell(z_i).
\]

Let $W \in \mathbb{R}^{k \times m}$ be a weight matrix,
and $X \in \mathbb{R}^{n \times m}$ denote the observed data matrix.
Conditioned on $\vect{z}{}$ and $W$, the entries of $X$ are independent Gaussian random variables:
\begin{equation*}
    p(X \mid \vect{z}{}, W)
    =
    \prod_{i=1}^n \prod_{j=1}^m \Nor{X_{ij}}{(ZW)_{ij}}{\beta^{-1}},
\end{equation*}
where $\beta$ is a known precision parameter.

Thus, conditioned on $\vect{z}{}$, the columns of $X$ are independent and each column
$\vect{x}{j}$ corresponds to a separate linear regression model with its own weight vector $\vect{w}{j}$.

We assume a noninformative prior over the latent variable parameters:
\[
    p(\mu,\tau) \propto \tau^{-1},
\]
and place independent Gaussian priors on the columns of $W$:
\begin{equation*}
    \vect{w}{j} \sim \mathcal{N}\big(0, \alpha_j^{-1} I \big),
    \qquad j = 1, \dots, m,
\end{equation*}
where $\alpha_j > 0$ is a known precision hyperparameter associated with column $j$.

In order to formulate a variational treatment of this model,
we consider a variational distribution which factorizes between the latent
variables and the parameters so that
\[
    q(\vect{z}{}, W, \mu, \tau) = q(\vect{z}{})q(W, \mu, \tau).
\]

First, let us consider the derivation of the update equation for the factor $q(W, \mu, \tau)$.
The log of the optimized factor is given by
\begin{equation*}
    \ln q^*(W, \mu, \tau)
    = \Ex{q(\vect{z}{})}{\ln p(W, \mu, \tau \mid \vect{z}{}, X)} + \text{const}.
\end{equation*}
As can be seen from Figure~\ref{fig:regression_bayes_net}, the d-separation criterion implies that
\begin{equation*}
    \Ex{q(\vect{z}{})}{\ln p(W, \mu, \tau \mid \vect{z}{}, X)} =
    \Ex{q(\vect{z}{})}{\ln p(\mu, \tau \mid \vect{z}{})} +
    \sum_j \Ex{q(\vect{z}{})}{\ln p(\vect{w}{j} \mid \vect{z}{}, X)}.
\end{equation*}
Therefore, we have the following induced factorization for $q(W, \mu, \tau)$:
\begin{equation*}
    q(W, \mu, \tau) =  q(\mu, \tau) \prod_{j=1}^m q(\vect{w}{j}).
\end{equation*}

We obtain $q(\vect{w}{j})$ using the fact that each $\vect{w}{j}$ is a weight vector of a bayesian linear regression
model:
\begin{align*}
    \ln p(\vect{w}{j} \mid \vect{z}{}, X) &= -\hf[\beta] \norm{Z\vect{w}{j} - \vect{x}{j}} -
    \hf[\alpha_j] \trans{\vect{w}{j}} \vect{w}{j} + \text{const} \\
    &= -\hf \trans{\vect{w}{j}} (\beta \trans{Z}Z + \alpha_j I) \vect{w}{j} +
    \beta \trans{\vect{x}{j}} Z \vect{w}{j} + \text{const}.
\end{align*}
Taking the expectation with respect to $q(\vect{z}{})$, we obtain
\begin{equation*}
    \ln q^*(\vect{w}{j}) = -\hf \trans{\vect{w}{j}} (\beta \Ex{q(\vect{z}{})}{\trans{Z}Z} + \alpha_j I)
    \vect{w}{j} +
    \beta \trans{\vect{x}{j}} \Ex{q(\vect{z}{})}{Z} \vect{w}{j} + \text{const}.
\end{equation*}
The expression above is quadratic in $\vect{w}{j}$.
Since the logarithm of a multivariate Gaussian distribution is a quadratic
function of its argument, it follows that
\begin{equation*}
    q^*(\vect{w}{j}) = \mathcal{N}(\vect{w}{j} \mid \mu_j, \Sigma_j),
\end{equation*}
We therefore have
\begin{align*}
    \Ex{}{\vect{w}{j}} &= \mu_j, \\
    \Ex{}{\vect{w}{j} \trans{\vect{w}{j}}} &= \Sigma_j + \mu_j \trans{\mu_j},
\end{align*}
where
\begin{align*}
    \Sigma_j^{-1} &= \beta \Ex{q(\vect{z}{})}{\trans{Z}Z} + \alpha_j I, \\
    \mu_j &= \beta \Sigma_j \trans{\Ex{q(\vect{z}{})}{Z}} \vect{x}{j}.
\end{align*}

We now consider the update equation for the factor $q(\mu, \tau)$:
\begin{align*}
    \ln q^*(\mu, \tau)
    &= \Ex{q(\vect{z}{})}{ \ln p(\mu, \tau \mid \vect{z}{}) } + \text{const} \\
    &= - \frac{\tau}{2} \left(\sum_{i} \Ex{q(z_i)}{z_i^2} - 2\mu \sum_{i} \Ex{q(z_i)}{z_i} + n\mu^2 \right) \\
    &\quad + \left(\frac{n}{2}-1\right)\ln\tau  + \text{const}.
\end{align*}

We define
\begin{align*}
    \avgz &= \frac{1}{n} \sum_i \Ex{q(z_i)}{z_i}, \\
    \avgzz &= \frac{1}{n} \sum_i \Ex{q(z_i)}{z_i^2}.
\end{align*}
We observe that $q^*(\mu, \tau)$, for $n > 1$ follows a Normal–Inverse–Gamma distribution:
\begin{equation*}
    q^*(\mu, \tau) = \Nor{\mu}{\avgz}{\frac{1}{n \tau}} \Gam{\tau}{\hf[n-1]}{\hf[n(\avgzz - \avgz^{2})]}.
\end{equation*}
This implies
\begin{align*}
    \Ex{q(\tau)}{\tau} &= \frac{(n - 1)}{n(\avgzz - \avgz^{2})}, \\
    \Ex{q(\mu, \tau)}{\mu \tau} &= \avgz \Ex{q(\tau)}{\tau} \cdot
\end{align*}

Finally, we derive the update equation for $q(\vect{z}{})$:
\begin{align*}
    \ln q^*(\vect{z}{}) &= \Ex{q(W,\mu,\tau)}{\ln p(X \mid \vect{z}{}, W) +
    \ln p(\vect{z}{} \mid \mu, \tau)} + \text{const} \\
    &= \Ex{q(W)}{\ln p(X \mid \vect{z}{}, W)} +
    \Ex{q(\mu,\tau)}{\ln p(\vect{z}{} \mid \mu, \tau)} + \text{const}.
\end{align*}

We expand the terms $\ln p(X \mid \vect{z}{}, W)$ and $\ln p(\vect{z}{} \mid \mu, \tau)$.
All expressions that are independent of $\vect{z}{}$ are treated as constants.
\begin{align*}
    \ln p(X \mid \vect{z}{}, W)
    &= \sum_{j} \ln \mathcal{N} (\vect{x}{j} \mid Z \vect{w}{j},  \beta^{-1}I) \\
    &= -\hf[\beta] \sum_j \norm{Z\vect{w}{j} - \vect{x}{j}} + \text{const} \\
    &= \beta \sum_{i} \left(-\hf \trans{\phi (z_i)} \sum_j \vect{w}{j} \trans{\vect{w}{j}} \phi (z_i)  +
                          \trans{\phi(z_i)} \sum_j X_{ij} \vect{w}{j} \right) + \text{const}. \\ \\
    \ln p(\vect{z}{} \mid \mu, \tau) &= -\hf[\tau] \sum_i (z_i - \mu) ^ 2 + \text{const} \\
    &= \sum_i \left( -\hf[\tau] z_i^2 + \tau \mu z_i \right) + \text{const}.
\end{align*}
Substituting these expansions into the variational update expression yields:
\begin{equation*}
    \ln q^*(\vect{z}{}) = \sum_i \ln q^*(z_i),
\end{equation*}
where
\begin{align*}
    \ln q^*(z_i) &=
    \beta \trans{\phi (z_i)}
    \sum_{i} \left(-\hf \sum_j \Ex{q(\vect{w}{j})}{\vect{w}{j} \trans{\vect{w}{j}}} \phi (z_i) +
                 \sum_j X_{ij} \Ex{q(\vect{w}{j})}{\vect{w}{j}} \right) \\
    &\quad + \left( -\hf \Ex{q(\tau)}{\tau} z_i^2 + \Ex{q(\mu, \tau)}{\tau \mu} z_i \right)
    + \text{const}.
\end{align*}

This implies that $q(\vect{z}{})$ factorizes over its components $z_i$.
The update equations therefore require expectations of the form
\begin{equation*}
    \Ex{q(z_k)}{\phi_i(z_k)\phi_j(z_k)}, \quad
    \Ex{q(z_k)}{\phi_i(z_k)}, \quad
    \Ex{q(z_k)}{z_k}, \quad
    \Ex{q(z_k)}{z_k^2},
\end{equation*}
for all relevant indices $i,j,k$.
These expectations can be evaluated numerically.